\chapter{Pendahuluan}

\section{Latar Belakang}

Transformasi digital pada sektor sistem pembayaran Indonesia berlangsung
sangat cepat dalam satu dekade terakhir. Penetrasi telepon pintar dan
pertumbuhan ekonomi digital mendorong masyarakat beralih dari uang tunai menuju
instrumen pembayaran nontunai. Bank Indonesia melalui Blueprint Sistem
Pembayaran Indonesia (BSPI) menempatkan digitalisasi sebagai agenda utama untuk
menjaga kedaulatan moneter, mendorong inklusi keuangan, serta memperkuat
integrasi ekonomi dan keuangan digital nasional \cite{bi_bspi2025}. Dalam
konteks tersebut, kehadiran mata uang digital bank sentral atau \textit{Central
Bank Digital Currency} (CBDC) menjadi salah satu inisiatif strategis untuk
melengkapi uang kartal dan uang elektronik yang telah ada.

Pertumbuhan ekosistem digital tersebut, bagaimanapun, juga membawa risiko
baru bagi stabilitas sistem keuangan. Pada tahun 2022, penetrasi internet
Indonesia mencapai 74\% dengan 204,7 juta pengguna, sementara pelanggan
seluler mencapai 370,1 juta atau setara 133\% populasi. Pandemi COVID-19
menambah 21 juta konsumen digital baru yang sebelumnya belum terlayani
secara daring, sebagaimana dirangkum dalam Tabel 1 White Paper CBDC Bank
Indonesia \cite{bi_whitepaper2022}. Pada periode yang sama, kapitalisasi pasar
aset kripto global melonjak hingga hampir 3 triliun dolar AS, tumbuh
1.393\% antara 2019 dan 2021, sementara investor aset kripto di Indonesia
mencapai 16,3 juta per September 2022 dengan laju pertumbuhan 81,6\%
\textit{year-over-year} \cite{bi_whitepaper2022}. Fenomena ini memunculkan
\textit{cryptoization}, yaitu kondisi di mana adopsi aset kripto oleh
penduduk menyebar lebih cepat dibandingkan mata uang resmi
\cite{bi_whitepaper2022}. Disrupsi digital tidak lagi terbatas
pada isu \textit{shadow banking}, melainkan telah merambah pada fenomena
\textit{shadow currency} dan \textit{shadow central banking} yang mengikis
efektivitas kebijakan bank sentral \cite{bi_whitepaper2022}.
\textit{Stablecoin} dan aset kripto tanpa aset dasar memperbesar risiko
makrofinansial berupa melemahnya kendali otoritas atas stabilitas moneter dan
sistem keuangan \cite{bi_whitepaper2022}. Mata uang digital bank
sentral dirancang untuk menjaga kepercayaan publik terhadap uang dan
kedaulatan moneter, sekaligus mempertahankan peran bank sentral sebagai
jangkar stabilitas dalam ekonomi digital
\cite{bi_whitepaper2022, Auer2022Central}.

CBDC merupakan representasi digital dari uang yang diterbitkan dan menjadi
kewajiban langsung bank sentral, sehingga memiliki kualitas bebas risiko yang
setara dengan uang tunai \cite{unknown2024p}. Berbeda dengan uang elektronik
yang diterbitkan oleh penyelenggara swasta dan merupakan kewajiban penerbitnya,
CBDC menempatkan bank sentral sebagai penjamin nilai akhir. Bagi Indonesia,
Bank Indonesia menamai inisiatif CBDC sebagai Rupiah Digital melalui Proyek
Garuda, yang menempuh pendekatan bertahap dimulai dari Rupiah Digital
\textit{wholesale} (w-Digital Rupiah) untuk transaksi antarlembaga, kemudian
diperluas menuju Rupiah Digital \textit{retail} (r-Digital Rupiah) bagi
masyarakat luas \cite{bi_whitepaper2022, bi_garuda_poc}. Desain Rupiah Digital
mengadopsi model dua jenjang (\textit{two-tier}), yaitu Bank Indonesia
menerbitkan dan mengelola Rupiah Digital pada tingkat \textit{wholesale},
sementara distribusi kepada masyarakat dilakukan melalui pihak perantara seperti
perbankan dan penyedia jasa pembayaran \cite{bi_consultative}.

Salah satu pilihan teknologi yang banyak dikaji untuk menopang CBDC adalah
\textit{Distributed Ledger Technology} (DLT) berbasis \textit{blockchain}.
Teknologi ini menawarkan keunggulan berupa integritas data yang terjaga melalui
catatan yang bersifat \textit{append-only}, jejak audit yang
transparan, serta kemampuan otomasi aturan bisnis melalui \textit{smart
	contract} \cite{unknown2022, sethaput2025}. Untuk kebutuhan bank sentral, jenis
\textit{blockchain} yang sesuai adalah \textit{permissioned blockchain}, yaitu
jaringan dengan keanggotaan terbatas dan teridentifikasi, sehingga bank sentral
tetap memegang kendali atas tata kelola jaringan \cite{bhawana2021}. Hyperledger
Fabric merupakan salah satu kerangka kerja \textit{permissioned blockchain}
paling matang yang banyak digunakan pada ranah keuangan. Fabric mendukung
arsitektur modular, identitas berbasis \textit{Membership Service Provider}
(MSP), pemisahan transaksi melalui kanal (\textit{channel}), serta model
eksekusi \textit{execute-order-validate} yang mendukung pemrosesan paralel
tanpa bergantung pada mekanisme \textit{proof-of-work} yang boros energi
\cite{androulaki2018}.

Meskipun kajian mengenai CBDC berkembang pesat, fokus literaturnya masih
timpang. Sebagian besar penelitian menelaah CBDC dari sisi makroekonomi,
kebijakan moneter, dan stabilitas keuangan \cite{Barrdear2021The,
	Auer2022Central, unknown2024p}, atau dari sisi taksonomi desain tingkat tinggi
seperti pilihan \textit{token-based} versus \textit{account-based} dan model
satu jenjang versus dua jenjang \cite{Auer2020The, Kiff2020A, Allen2020Design,
	Tsareva2024Retail}. Sebaliknya, artefak teknis yang utuh, terbuka, dan dapat
diuji secara empiris untuk konteks CBDC ritel masih jarang tersedia.

Kajian teknis yang ditinjau masih membahas komponen CBDC secara terpisah.
Privasi dan auditabilitas dibahas dalam \cite{islam2023, Pocher2022Privacy,
Michalopoulos2025Privacy}, keamanan dan pemodelan ancaman dalam
\cite{hans2023}, serta arsitektur jaringan dalam \cite{bhawana2021, sun2017}.
Kelayakan sistem CBDC ritel, sebaliknya, bergantung pada keterpaduan komponen
tersebut. Dalam korpus literatur yang digunakan pada penelitian ini, penerbitan
dan distribusi likuiditas berjenjang, transfer ritel, serta pengawasan belum
dievaluasi sebagai satu artefak Fabric yang kinerjanya dapat diukur.

Untuk konteks Indonesia, Bank Indonesia telah menetapkan arah desain Rupiah
Digital berupa model dua jenjang dan interoperabilitas sistem pembayaran
\cite{bi_whitepaper2022, bi_consultative}. Desain CBDC bagi perekonomian
Indonesia juga telah dirumuskan dalam \cite{zams2020delphi,
	Syarifuddin2024OPTIMAL, aryani2023critlit}. Namun, realisasi teknis Proyek
Garuda yang terdokumentasi secara publik baru mencakup lapisan
\textit{wholesale} \cite{bi_garuda_poc}, sedangkan purwarupa lapisan
\textit{retail} yang selaras dengan spesifikasi tersebut belum tersedia secara
terbuka. Dari sisi evaluasi kinerja, CBDC ritel yang telah diluncurkan maupun
yang masih berupa pilot (seperti eNaira di Nigeria \cite{ree2023enaira} dan
pilot ritel di Thailand \cite{Bhensook2025Public}) belum menyertakan tolok ukur
kinerja terstandar yang dapat dibandingkan langsung dengan purwarupa ini. Sistem
berkinerja tinggi seperti
Hamilton/OpenCBDC menggunakan arsitektur \textit{account-based} terpusat,
bukan model \textit{blockchain} \cite{Lovejoy2022A}.
Akibatnya, data \textit{throughput} dan latensi yang terukur untuk skenario CBDC
ritel di atas \textit{permissioned blockchain}, khususnya untuk Rupiah Digital,
masih langka.

Atas dasar literatur dan artefak publik yang ditinjau, celah penelitian berada
pada kebutuhan akan purwarupa CBDC ritel yang (i) mengintegrasikan penerbitan
berjenjang, transfer ritel, serta pengawasan dalam satu sistem, (ii) selaras
dengan rancangan dua jenjang Bank Indonesia, dan (iii) dievaluasi melalui
pengujian Hyperledger Caliper yang dapat diulang dengan konfigurasi dan artefak
yang dicatat. Penelitian ini
merancang dan
membangun sistem Rupiah Digital ritel berbasis Hyperledger Fabric yang
merepresentasikan ekosistem multipihak: Bank Indonesia, bank validator,
penyedia jasa pembayaran (PJP), pedagang, pelanggan ritel, pengawas, serta
pengamat dari Otoritas Jasa Keuangan (OJK). Sistem tersebut dilengkapi
jalur baca pengawasan dan catatan audit pada buku besar. Evaluasi dengan Hyperledger Caliper
mengukur \textit{throughput} dan latensi pada berbagai skenario beban untuk
menilai karakteristik kinerjanya dalam ruang lingkup purwarupa.

\section{Rumusan Masalah}

Berdasarkan latar belakang di atas, rumusan masalah pada penelitian ini adalah
sebagai berikut.
\begin{enumerate}
	\item Bagaimana merancang dan membangun sistem Rupiah Digital ritel berbasis
	      Hyperledger Fabric yang mendukung penerbitan oleh Bank Indonesia ke dalam
	      Treasury, distribusi langsung kepada bank validator atau PJP sebagai
	      kustodian, serta layanan kepada pelanggan ritel, transfer ritel, dan
	      pengawasan transaksi dengan batas
	      wewenang organisasi yang jelas?
	\item Bagaimana kinerja sistem yang dibangun dari sisi \textit{throughput}
	      dan latensi pada berbagai skenario beban transaksi yang diukur menggunakan
	      Hyperledger Caliper?
\end{enumerate}

\section{Tujuan Penelitian}

Tujuan yang ingin dicapai pada penelitian ini adalah sebagai berikut.
\begin{enumerate}
	\item Merancang dan membangun purwarupa sistem Rupiah Digital ritel berbasis
	      Hyperledger Fabric dengan arsitektur multipihak yang mendukung penerbitan
	      ke Treasury Bank Indonesia, distribusi langsung kepada bank validator atau
	      PJP sebagai kustodian, transfer ritel, dan pengawasan transaksi dengan batas
	      wewenang organisasi yang jelas.
	\item Mengevaluasi kinerja sistem dari sisi \textit{throughput} dan latensi
	      pada berbagai skenario beban menggunakan Hyperledger Caliper.
\end{enumerate}

\section{Batasan Penelitian}

\begin{enumerate}
	\item Sistem yang dibangun merupakan purwarupa (\textit{prototype}) yang
	      dijalankan pada lingkungan jaringan tersimulasi secara lokal menggunakan
	      \textit{container}, bukan sistem produksi.
	\item Jaringan \textit{blockchain} menggunakan Hyperledger Fabric versi 2.5
	      dengan lima organisasi peserta dan mekanisme konsensus \textit{ordering}
	      berbasis Raft (etcdraft).
	\item Penelitian berfokus pada perancangan, implementasi, dan evaluasi
	      kinerja fungsional; aspek hukum, kebijakan moneter, dan dampak makroekonomi
	      dari penerapan Rupiah Digital berada di luar cakupan. Analisis strategi
	      kebijakan moneter dan dampak makroekonomi CBDC, termasuk peran
	      \textit{CBDC rate}, didelegasikan pada studi terpisah seperti
	      \cite{SyarifuddinBakhtiar2021Monetary}.
	\item Evaluasi kinerja dilakukan menggunakan Hyperledger Caliper dengan beban
	      yang dihasilkan oleh sejumlah \textit{worker} pada lingkungan pengujian lokal.
	\item Pengujian dijalankan pada satu inang
	      (\textit{single host}) dengan satu, dua, dan empat \textit{worker} Caliper dan satu
	      \textit{peer} per organisasi, tanpa latensi jaringan dan persebaran geografis
	      yang nyata. Karena itu, angka \textit{throughput} dan latensi yang diperoleh
	      tidak dapat langsung digeneralisasi ke perangkat keras maupun topologi
	      produksi \cite{Lovejoy2022A, fabric_perf2023}.
	\item Komposisi beban transaksi dan distribusi
	      nilai yang digunakan merupakan asumsi desain yang menyerupai pola pembayaran
	      ritel di Indonesia, bukan jejak transaksi empiris. Profil transfer menggunakan
	      konfigurasi terpisah untuk satu, dua, dan empat \textit{worker}; setiap profil
	      menetapkan populasi, slot transaksi, saldo awal, dan durasi pada berkas
	      konfigurasi yang ikut dihash dalam manifes. Ronde pemanasan berlangsung
	      30 detik, sedangkan ronde terukur berlangsung 120 detik.
	\item KYC dan QRIS tersedia sebagai fitur implementasi pendukung, tetapi tidak
	      menjadi kriteria penerimaan utama penelitian.
	\item Evaluasi transfer dasar menggunakan skenario perebutan kunci rendah
	      (\textit{low-contention}) dengan rencana transaksi berputar dan populasi
	      tetap. Skenario terpisah menguji perebutan satu dompet melalui permintaan
	      transfer agregat yang melebihi saldo awal. Penelitian tidak menggunakan
	      jejak transaksi empiris untuk menetapkan tingkat konsentrasi maupun
	      kebijakan \textit{retry/backoff} yang representatif. Hasilnya karena itu
	      tidak mewakili beban pedagang populer atau kapasitas produksi.
	\item Purwarupa ini tidak
	      mengimplementasikan pembayaran CBDC tanpa koneksi jaringan. Kemampuan
	      tersebut memerlukan desain perangkat aman, verifikasi kriptografis,
	      pencegahan \textit{double-spending}, dan rekonsiliasi setelah konektivitas
	      kembali \cite{polaris_security, polaris_handbook}. Hasil penelitian ini
	      hanya merepresentasikan transaksi yang diselesaikan melalui
	      koneksi ke \textit{ledger}.
	\item Sistem belum menyediakan privasi transaksi pada kanal bersama.
	      Topologi memakai satu \textit{application channel} tanpa \textit{private
	      data collection} atau teknik kriptografi penjaga privasi. Akibatnya,
	      metadata dompet dan graf transaksi terlihat oleh peserta kanal yang
	      berwenang. Purwarupa ini belum menangani kompromi antara privasi dan keterlacakan
	      \cite{Pocher2022Privacy, Soana2024Central, Said2026Designing}.
	\item Topologi hanya memiliki satu \textit{consenter} Raft sehingga tidak
	      menoleransi kegagalan \textit{orderer}. Evaluasi juga tidak mencakup
	      kompromi simpul, ketersediaan tinggi, RTO/RPO, rotasi atau revokasi
	      identitas, peningkatan versi, respons insiden, dan pemulihan bencana.
	      Kesiapan produksi serta dampak operasional dan inklusi keuangan tidak
	      dinilai.
	\item Proyek Garuda menetapkan tiga prinsip 3i sebagai prasyarat
	      desain Rupiah Digital: integrasi dengan infrastruktur pasar keuangan (IPK)
	      domestik, interoperabilitas teknis dan semantik dengan sistem yang ada maupun
	      yang akan datang, serta interkoneksi lintas batas untuk jaringan CBDC
	      antarnegara \cite{bi_whitepaper2022}. Purwarupa ini memenuhi prinsip tersebut
	      hanya pada tataran internal: alur penerbitan--distribusi--pembayaran--pengawasan
	      berjalan secara terpadu dalam satu jaringan Hyperledger Fabric. Namun,
	      integrasi eksternal dengan BI-RTGS dan BI-FAST hanya dimodelkan sebagai titik
	      masuk (\textit{stub}) tanpa koneksi nyata, Penyimpan Digital Rupiah (KDR)
	      belum diimplementasikan, konformitas terhadap standar SNAP API nasional tidak
	      diterapkan, dan tidak ada gerbang (\textit{gateway}) menuju IPK domestik
	      maupun jaringan CBDC lintas batas. Pemenuhan prinsip 3i secara penuh
	      merupakan pekerjaan lanjutan yang berada di luar cakupan
	      purwarupa ini.
\end{enumerate}

\section{Manfaat Penelitian}

Penelitian ini diharapkan memberikan manfaat sebagai berikut.
\begin{enumerate}
	\item Memberikan rujukan teoretis tentang perancangan CBDC ritel berbasis
	      \textit{permissioned blockchain} yang mencakup penerbitan berjenjang,
	      distribusi kepada kustodian bank atau PJP, transaksi antarpengguna, dan pengawasan oleh
	      bank sentral.
	\item Menjelaskan penerapan Hyperledger Fabric untuk memodelkan pemisahan
	      peran antara BI, peserta distribusi, dan pengguna ritel tanpa
	      menempatkan data portal pada buku besar.
	\item Menyediakan purwarupa yang dapat digunakan sebagai bahan eksperimen
	      teknis untuk menguji alur penerbitan, transfer, otorisasi, dan
	      pencatatan peristiwa pengawasan.
	\item Memberikan dasar evaluasi kinerja awal melalui pengujian beban
	      terbatas sehingga karakteristik \textit{throughput}, latensi, dan
	      keberhasilan transaksi dapat dibahas secara terukur dalam ruang
	      lingkup purwarupa.
	\item Mendukung diskusi teknis pengembangan Rupiah Digital ritel dengan menunjukkan
	      batas kemampuan purwarupa, termasuk kebutuhan perluasan topologi,
	      privasi transaksi, dan pengujian keamanan lanjutan, tanpa
	      menyimpulkan dampak inklusi atau kesiapan produksi.
\end{enumerate}

\section{Sistematika Penulisan}

Penulisan skripsi ini disusun ke dalam lima bab dengan sistematika sebagai
berikut.

\noindent BAB I PENDAHULUAN berisi latar belakang, rumusan masalah,
tujuan penelitian, batasan penelitian, manfaat penelitian, dan sistematika
penulisan.

\noindent BAB II TINJAUAN PUSTAKA DAN DASAR TEORI memaparkan tinjauan
terhadap penelitian terdahulu yang relevan serta dasar teori yang mendasari
penelitian, meliputi CBDC, \textit{blockchain} dan DLT, Hyperledger Fabric,
dan Hyperledger Caliper.

\noindent BAB III METODE PENELITIAN menguraikan alat dan bahan, metode
perancangan, metode implementasi, metode pengujian, serta alur penelitian yang
ditempuh.

\noindent BAB IV HASIL DAN PEMBAHASAN menyajikan hasil implementasi
sistem, hasil evaluasi kinerja menggunakan Hyperledger Caliper, serta
perbandingan dengan penelitian terdahulu.

\noindent BAB V KESIMPULAN DAN SARAN memuat kesimpulan yang menjawab
rumusan masalah serta saran untuk pengembangan penelitian selanjutnya.
